% Chapter 3: Introduction to multiparameter models
% Bayesian Data Analysis - Gelman et al.

\chapter{多参数模型简介}

\section{对 Nuisance 参数的积分}

为了表达联合后验分布和边际后验分布的思想，假设 $\theta$ 有两部分，每部分都可以是向量：$\theta = (\theta_1, \theta_2)$。进一步假设我们只（至少目前只）对 $\theta_1$ 的推断感兴趣，因此 $\theta_2$ 可以被视为 nuisance 参数。

例如，在简单例子 $y|\mu, \sigma^2 \sim \text{Normal}(\mu, \sigma^2)$ 中，$\mu = \theta_1$ 和 $\sigma^2 = \theta_2$ 都未知，但通常我们最关心的是 $\mu$。

我们寻求给定观测数据的参数条件分布，即 $p(\theta_1|y)$。这从联合后验密度
\[
p(\theta_1, \theta_2|y) \propto p(y|\theta_1, \theta_2) \times p(\theta_1, \theta_2)
\]
通过对 $\theta_2$ 积分得到：
\[
p(\theta_1|y) = \int p(\theta_1, \theta_2|y) \, d\theta_2
\]

另一种分解方式是：
\[
p(\theta_1|y) = \int p(\theta_1|\theta_2, y) \times p(\theta_2|y) \, d\theta_2
\]

它表明：感兴趣的\textbf{后验分布是条件后验分布的混合}，其中 $p(\theta_2|y)$ 是对 $\theta_2$ 不同值的加权函数。权重依赖于 $\theta_2$ 的后验密度，因此依赖于数据和先验模型的组合。

对 nuisance 参数 $\theta_2$ 的积分可以一般性地解释。例如，$\theta_2$ 可以包括代表不同可能子模型的离散成分。

\begin{Remark}
  我们很少直接计算这个积分，但它暗示了构建和计算多参数模型的重要策略：\textbf{后验分布可以通过边际和条件模拟来计算}——首先从边际后验分布抽取 $\theta_2$，然后在给定 $\theta_2$ 值的条件下从条件后验分布抽取 $\theta_1$。通过这种方式，积分被间接执行。
\end{Remark}

这种分析形式的典型例子是正态模型中均值和方差都未知的情况。

\section{正态数据与非信息先验}

作为从样本估计总体均值的原型例子，我们考虑来自一元正态分布 $\text{Normal}(\mu, \sigma^2)$ 的 $n$ 个独立观测 $y = (y_1, \ldots, y_n)$。

我们首先在非信息先验分布下分析这个模型。

\subsection{非信息先验分布}

对于 $\mu$ 和 $\sigma^2$，一个合理的模糊先验密度，假设位置参数和尺度参数的先验独立，是在 $(\mu, \log \sigma^2)$ 上均匀，即：
\[
p(\mu, \sigma^2) \propto (\sigma^2)^{-1}
\]

\subsection{联合后验分布}

在这种约定的非正常先验密度下，联合后验分布为：
\[
p(\mu, \sigma^2 | y) \propto (\sigma^2)^{-(n/2+1)} \exp\left(-\frac{1}{2\sigma^2} \sum_{i=1}^n (y_i - \mu)^2\right)
\]

利用 $\sum_{i=1}^n (y_i - \mu)^2 = \sum_{i=1}^n (y_i - \bar{y})^2 + n(\bar{y} - \mu)^2$，可以写成：
\[
= (\sigma^2)^{-(n/2+1)} \exp\left(-\frac{1}{2\sigma^2}\left[(n-1)s^2 + n(\bar{y} - \mu)^2\right]\right)
\]
其中 $\bar{y}$ 是样本均值，$s^2 = \frac{1}{n-1}\sum_{i=1}^n (y_i - \bar{y})^2$ 是样本方差。

\textbf{充分统计量}是 $\bar{y}$ 和 $s^2$。

\subsection{条件后验分布}

为了分解联合后验密度，我们首先考虑条件后验密度 $p(\mu | \sigma^2, y)$，然后考虑边际后验密度 $p(\sigma^2 | y)$。

对于给定 $\sigma^2$ 的 $\mu$ 的后验分布，直接使用第二章中已知方差正态均值的均匀先验结果：
\[
\mu | \sigma^2, y \sim N(\bar{y}, \sigma^2 / n)
\]

\subsection{边际后验分布}

为了确定 $p(\sigma^2 |y)$，我们必须对 $\mu$ 积分：
\[
p(\sigma^2 | y) \propto \int (\sigma^2)^{-(n/2+1)} \exp\left(-\frac{1}{2\sigma^2}\left[(n-1)s^2 + n(\bar{y} - \mu)^2\right]\right) d\mu
\]

对 $\mu$ 积分需要计算简单的正态积分，结果是：
\[
p(\sigma^2 | y) \propto (\sigma^2)^{-(n/2+1)} \exp\left(-\frac{(n-1)s^2}{2\sigma^2}\right) \times \sqrt{2\pi\sigma^2/n}
\]

简化后：
\[
p(\sigma^2 | y) \propto (\sigma^2)^{-[(n-1)/2+1]} \exp\left(-\frac{(n-1)s^2}{2\sigma^2}\right)
\]

这是尺度逆卡方分布（scale inverse-chi-squared distribution）：

\begin{Definition}[尺度逆卡方分布]
  \[
  p(x|\nu, \sigma^2) = \frac{(\nu\sigma^2/2)^{\nu/2}}{\Gamma(\nu/2)} x^{-(\nu/2+1)} \exp\left(-\frac{\nu\sigma^2}{2x}\right)
  \]
  其中 $\nu$ 是自由度，$\sigma^2$ 是尺度参数。
\end{Definition}

因此：
\[
\sigma^2 | y \sim \text{Inv-}\chi^2(n-1, s^2)
\]

由 $\frac{(n-1)s^2}{\sigma^2} \sim \chi_{n-1}^2$，这与频率统计中的经典结果有显著相似性。

\subsection{从联合后验分布抽样}

从联合后验分布抽取样本很容易：
\begin{enumerate}
  \item 从边际后验分布抽取 $\sigma^2$，即 $\sigma^2 \sim \text{Inv-}\chi^2(n-1, s^2)$
  \item 从条件后验分布抽取 $\mu$，即 $\mu | \sigma^2, y \sim N(\bar{y}, \sigma^2/n)$
\end{enumerate}

\subsection{边际后验分布：$t$ 分布}

总体均值 $\mu$ 通常是感兴趣的估计量，边际后验分布可以通过对 $\sigma^2$ 积分得到。结果是：
\[
\mu | y \sim t_{n-1}(\bar{y}, s^2/n)
\]
即自由度为 $n-1$ 的学生 $t$ 分布。

这与频率统计中的经典结果也有显著相似性：在采样分布下，$\frac{\bar{y}-\mu}{s/\sqrt{n}} \sim t_{n-1}$。

一个值得注意的事实是：这个边际后验分布不依赖于数据——它是数据的函数，但关于 $\mu$ 的概率陈述只通过 $\bar{y}$ 和 $s^2$ 依赖于数据。

\subsection{后验预测分布}

未来观测 $\tilde{y}$ 的后验预测分布可以写成混合：
\[
p(\tilde{y} | y) = \iint p(\tilde{y}|\mu, \sigma^2) \times p(\mu, \sigma^2 | y) \, d\mu \, d\sigma^2
\]

结果是 $t$ 分布：$\tilde{y} \sim t_{n-1}(\bar{y}, (1+1/n)s^2)$。

\section{正态数据与共轭先验}

当使用共轭先验时，分析更加系统化。正态模型的共轭先验族具有形式：
\[
\mu | \sigma^2 \sim N(\mu_0, \sigma^2/\kappa_0), \quad \sigma^2 \sim \text{Inv-}\chi^2(\nu_0, \sigma_0^2)
\]

这对应于联合先验密度 $N\text{-Inv-}\chi^2(\mu_0, \sigma_0^2/\kappa_0; \nu_0, \sigma_0^2)$。

四个超参数分别是 $\mu$ 的位置和尺度、$\sigma^2$ 的自由度和尺度。

注意 $\sigma^2$ 出现在 $\mu|\sigma^2$ 的条件分布中，意味着在联合共轭先验密度中 $\mu$ 和 $\sigma^2$ 必然相关。

这通常是有意义的——让 $\mu$ 的先验方差与 $\sigma^2$（观测 $y$ 的采样方差）挂钩是合理的。

\subsection{联合后验分布}

乘以正态似然得到后验密度仍然是 $N\text{-Inv-}\chi^2$ 形式，参数更新为：
\[
\mu_n = \frac{\kappa_0}{\kappa_0+n}\mu_0 + \frac{n}{\kappa_0+n}\bar{y}, \quad
\kappa_n = \kappa_0 + n, \quad
\nu_n = \nu_0 + n
\]
\[
\nu_n\sigma_n^2 = \nu_0\sigma_0^2 + (n-1)s^2 + \frac{\kappa_0 n}{\kappa_n}(\bar{y}-\mu_0)^2
\]

条件后验分布：$\mu | \sigma^2, y \sim N(\mu_n, \sigma^2/\kappa_n)$

边际后验分布：$\sigma^2 | y \sim \text{Inv-}\chi^2(\nu_n, \sigma_n^2)$

$\mu$ 的边际后验分布：$\mu | y \sim t_{\nu_n}(\mu_n, \sigma_n^2/\kappa_n)$

\section{分类数据的多项模型}

多项模型用于描述每个观测是 $k$ 个可能结果之一的数据。如果 $y$ 是每个结果观测次数的计数向量，则：
\[
p(y|\theta) \propto \prod_{j=1}^k \theta_j^{y_j}, \quad \sum_{j=1}^k \theta_j = 1
\]

共轭先验分布是 Dirichlet 分布的多元推广：
\[
p(\theta|\alpha) \propto \prod_{j=1}^k \theta_j^{\alpha_j-1}
\]

结果是 $\theta_j' s$ 的后验分布是 Dirichlet 分布，参数为 $\alpha_j + y_j$。

设置 $\alpha_j = 1$ 给出均匀密度；设置 $\alpha_j = 0$ 给出在 $\log(\theta_j)$ 上均匀的不正常先验。

\section{已知方差的多元正态模型}

多元正态模型 $N(\mu, \Sigma)$ 的似然函数为：
\[
p(y|\mu, \Sigma) \propto |\Sigma|^{-1/2} \exp\left\{-\frac{1}{2}(y-\mu)^T\Sigma^{-1}(y-\mu)\right\}
\]

给定 $\Sigma$ 时，$\mu$ 的共轭先验是多元正态分布 $\mu \sim N(\mu_0, \Lambda_0)$。

后验分布：
\[
\mu_n = (\Lambda_0^{-1} + n\Sigma^{-1})^{-1}(\Lambda_0^{-1}\mu_0 + n\Sigma^{-1}\bar{y}), \quad
\Lambda_n^{-1} = \Lambda_0^{-1} + n\Sigma^{-1}
\]

后验均值是数据和先验均值的加权平均，权重由精度矩阵给出；后验精度是先验精度和数据精度之和。

\section{均值和方差都未知的多元正态模型}

$(\mu, \Sigma)$ 的共轭先验分布是 normal-inverse-Wishart 分布：
\[
\Sigma \sim \text{Inv-Wishart}_{\nu_0}(\Lambda_0^{-1}), \quad \mu | \Sigma \sim N(\mu_0, \Sigma/\kappa_0)
\]

后验分布具有相同形式，参数更新为：
\[
\mu_n = \frac{\kappa_0}{\kappa_0+n}\mu_0 + \frac{n}{\kappa_0+n}\bar{y}, \quad
\kappa_n = \kappa_0 + n, \quad
\nu_n = \nu_0 + n
\]
\[
\Lambda_n = \Lambda_0 + S + \frac{\kappa_0 n}{\kappa_0+n}(\bar{y}-\mu_0)(\bar{y}-\mu_0)^T
\]
其中 $S = \sum_{i=1}^n (y_i - \bar{y})(y_i - \bar{y})^T$。

$\mu$ 的边际后验分布是多元 $t$ 分布。

\section{基本建模与计算总结}

多参数模型不允许轻易计算后验分布这一问题并非主要实际障碍，原因有三：

\begin{enumerate}
  \item 当参数较少时，非共轭多参数模型的后验推断可以通过简单模拟方法获得。
  \item 复杂模型通常可以表示为层次化或条件化形式，为此有有效的计算策略。
  \item 我们经常可以应用后验分布的正态近似，因此正态模型的共轭结构在实践中可以发挥重要作用，远远超出其对显式正态采样模型的应用。
\end{enumerate}

贝叶斯建模与计算的基本步骤总结：

\begin{enumerate}
  \item 写出模型的似然部分 $p(y|\theta)$，忽略任何不含 $\theta$ 的因子。
  \item 写出后验密度 $p(\theta|y) \propto p(\theta) \times p(y|\theta)$
    \begin{itemize}
      \item 如果先验信息表述良好，将其包含在 $p(\theta)$ 中。
      \item 否则使用弱信息先验或暂时设 $p(\theta) \propto \text{常数}$。
    \end{itemize}
  \item 创建参数的粗略估计 $\hat{\theta}$，用作起点和下一步计算的比较。
  \item 从后验分布抽取模拟 $\theta^1, \ldots, \theta^S$。使用样本模拟计算 $\theta$ 任何感兴趣函数的后验密度。
  \item 如果有预测量 $\tilde{y}$，通过对每个 $\tilde{y}^s$ 从给定 $\theta^s$ 的采样分布 $p(\tilde{y}|\theta^s)$ 抽取来模拟 $\tilde{y}^1, \ldots, \tilde{y}^S$。
\end{enumerate}

对于非共轭模型，步骤 4 可能很困难。各种方法已被开发用于从复杂模型中抽取后验模拟。

\section{本章小结}

本章我们学习了：
\begin{enumerate}
  \item 如何对 nuisance 参数进行积分（边际化和条件化）
  \item 多参数模型的计算策略：边际和条件模拟
  \item 正态模型中均值和方差同时未知时的分析
  \item 尺度逆卡方分布与 $t$ 分布
  \item 非信息先验和共轭先验在多参数情况下的应用
  \item 多项模型与 Dirichlet 先验
  \item 多元正态模型与 Inverse-Wishart 先验
\end{enumerate}

\section{下节预告}

下一章我们将学习渐近性以及贝叶斯方法与其他统计方法的联系。

\endinput
